% SPDX-License-Identifier: Apache-2.0
% 02_by_hand.tex: Drawing a block diagram by hand
\section{Drawing a Block Diagram by Hand}
\label{sec:by-hand}

An architecture diagram shows what a design does, not how synthesis built
it. Someone has to decide what goes in it, so it is written by hand. In
blockpic, it is written as text. Listing~\ref{lst:blockpic} is the counter's
architecture.

\lstinputlisting[style=bp,firstline=6,caption={The counter's architecture,
\code{counter.blockpic}.},label={lst:blockpic},float=t]{counter.blockpic}

A \code{block} has a label, a kind, and its input and output pins. The kind
picks the symbol: \code{ADD} draws an adder, \code{MUX} a multiplexer and
\code{DFF} a flip-flop. A \code{name}, such as \code{+1}, is drawn in place
of the label. An \code{INPUT\_PIN} or
\code{OUTPUT\_PIN} block is a port of the design. A \code{wire} joins an
output pin to an input pin. A wire with a \code{thickness} is drawn as a
bus: a thick line, crossed by a short slash with its width beside it. A
wire can also name the type of the bus, as in \code{type: "nibble\_t"};
the mark then shows the type in place of the width. A \code{group} draws
a dashed box around blocks that belong together, here the counter itself.

\begin{figure}[t]
\centering
\includegraphics[width=\columnwidth]{architecture.pdf}
\caption{The counter's architecture, drawn by blockpic from
Listing~\ref{lst:blockpic}. The thick lines, each marked 4, are the 4-bit buses.}
\label{fig:architecture}
\end{figure}

Fig.~\ref{fig:architecture} is the result. Blockpic placed every block and
routed every wire; the file gives no coordinates. The style is that of the
figures in AMD's UG474 \cite{ug474}: a plain sans-serif face, thin lines
without arrowheads, and pin names such as \code{D} and \code{CK} inside the
box. The feedback from the register to the adder is routed around the
outside, and every 4-bit bus is drawn thick and marked with its width.

\begin{guideline}[Keep the diagram beside the code]
Put the \code{.blockpic} file in the same directory as the HDL it describes.
A change to the design then shows the diagram in the same review, and the
reviewer can ask for the two to agree.
\end{guideline}

\begin{tip}
Name blocks after the signals they make. Here the register is called
\code{count}, as in the Verilog, so a reader can match the figure to the
listing by name.
\end{tip}
